
\input{config/documento}

\title{Agenda 2030 - Avaliação do avanço no Brasil dos Objetivos de Desenvolvimento Sustentável 01 Erradicação da pobreza e 02 Fome zero e agricultura sustentável da Agenda 2030}
\def\autor{Oscar Sérgio Monteiro\\
Vinicius Dutra\\
Jhonatan \\
\texttt{oscar.monteiro@ufpr.br}}
\def\disciplina{SC329 - Sustentabilidade  Ambiental e Social, turma b}
\def\professor{Elcídio Henriques  Quiraque}
\begin{document}
\input{config/pretextual}
\chapter{Introdução}
\subsection{Objetivos de Desenvolvimento Sistentável e Agenda 2030}
Os Objetivos de Desenvolvimento Sustentável são a evolução dos Objetivos de Desenvolvimento do Milênio, que foram propostos na Cúpula do Milênio em 2000, Nova York. Reunião promovida pela ONU que teve representantes de 189 países. Esse documento propôs   o linhas de ação:\\
 • redução da pobreza extrema,                                                  
 • educação,                                                                    
 • igualdade de gênero,                                                         
 • saúde,                                                                       
 • sustentabilidade ambiental,                                                  
 • combate a doenças,                                                           
 • parceria global para o desenvolvimento.                                      
Os Objetivos de desenvolvimento foram revisados na Cúpula das Nações Unidas de 2015, Nova York. Essa reunião  observou que: a questão da sustentabilidade deveria ser melhor detalhada abordando questões de consumo e produção sustentável; a preocupação com as questões climáticas, que não foi um ODM, deveria receber destaque; a finitude e a dependência dos recursos hídricos deveria receber um destaque; o problema se sustentabilidade deve ser tratado de forma universal e não se restringe ao desafio nas nações em desenvolvimento.
   	     	"A nova Agenda se baseia nos Objetivos de Desenvolvimento do Milênio e pretende completar o que estes não alcançaram, particularmente em alcançar os mais vulneráveis. "
A cúpula de 2015 definiu os seguintes 17 Objetivos de Desenvolvimento Sustentável:
 1 – Erradicação da pobreza    
 2 – Fome zero e agricultura sustentável                                 
 3 – Saúde e bem-estar                                                   
 4 – Educação de qualidade                                               
 5 – Igualdade de gênero                                                 
 6 – Água potável e saneamento                                           
 7 – Energia acessível e limpa                                           
 8 – Trabalho decente e crescimento econômico                            
 9 – Indústria, inovação e infraestrutura                                
 10 – Redução das desigualdades                                          
 11 – Cidades e comunidades sustentáveis                                 
 12 – Consumo e produção responsáveis                                    
 13 – Ação contra a mudança global do clima                              
 14 – Vida na água                                                       
 15 – Vida terrestre                                                     
 16 – Paz, justiça e instituições eficazes                               
 17 – Parcerias e meios de implementação Se quiser, posso também colocar
 
O documento Agenda 2030, sem SER alarmista, chama a atenção para a necessidade premente de enfrentar os desafios que ameaça a humanidade:
"Nosso mundo hoje

14. Encontramo-nos num momento de enormes desafios para o desenvolvimento sustentável. Bilhões de cidadãos continuam a viver na pobreza e a eles é negada uma vida digna. Há crescentes desigualdades dentro dos e entre os países. Há enormes disparidades de oportunidades, riqueza e poder. A desigualdade de gênero continua a ser um desafio fundamental. O desemprego, particularmente entre os jovens, é uma grande preocupação. Ameaças globais de saúde, desastres naturais mais frequentes e intensos, conflitos em ascensão, o extremismo violento, o terrorismo e as crises humanitárias relacionadas e o deslocamento forçado de pessoas ameaçam reverter grande parte do progresso do desenvolvimento feito nas últimas décadas.

O esgotamento dos recursos naturais e os impactos negativos da degradação ambiental, incluindo a desertificação, secas, a degradação dos solos, a escassez de água doce e a perda de biodiversidade acrescentam e exacerbam a lista de desafios que a humanidade enfrenta. A mudança climática é um dos maiores desafios do nosso tempo e seus efeitos negativos minam a capacidade de todos os países de alcançar o desenvolvimento sustentável. Os aumentos na temperatura global, o aumento do nível do mar, a acidificação dos oceanos e outros impactos das mudanças climáticas estão afetando seriamente as zonas costeiras e os países costeiros de baixa altitude, incluindo muitos países menos desenvolvidos e os pequenos Estados insulares em desenvolvimento. A sobrevivência de muitas sociedades, bem como dos sistemas biológicos do planeta, está em risco."
 Por ser uma agenda propositiva, da enfase ao principio liberal das "Ralações Internacionais" , enfatizando que os desafios podem e devem ser tratados pela cooperação e parceria. Porem as reflexões sobre causas do não atingimento das metas das ODM não são apresentadas. Permanecendo o dilema dos Estados Nações que existem para a segurança e desenvolvimento  dos seus cidadãos, com restrição de legitimidade para abrir mão da sua soberania em pról da segurança e desenvolvimento de cidadãos de outros estados. Desta forma o texto reafirma o princípio da soberania dos estados.
"Reafirmamos que cada Estado tem, e exerce livremente, sua soberania plena e permanente sobre toda a sua riqueza, seus recursos naturais e sua atividade econômica. Vamos implementar a Agenda para o pleno benefício de todos, para a geração de hoje e para as gerações futuras. Ao fazê-lo, reafirmamos nosso compromisso com o direito internacional e enfatizamos que a Agenda deverá ser implementada de uma forma consistente com os direitos e obrigações dos Estados sob o direito internacional."

 
Erradicação da Pobreza.
A Agende 2030 definiu as sequintes metas para a erradicação da pobreza:
 • 1.1 até 2030, erradicar a pobreza extrema para todas as pessoas em todos os  
   lugares, atualmente medida como pessoas vivendo com menos de US$ 1,25 por    
   dia.                                                           
 • 1.2 até 2030, reduzir pelo menos à metade a proporção de homens, mulheres e  
   crianças, de todas as idades, que vivem na pobreza, em todas as suas         
   dimensões, de acordo com as definições nacionais.                            
 • 1.3 implementar, em nível nacional, medidas e sistemas de proteção social    
   apropriados, para todos, incluindo pisos, e até 2030 atingir a cobertura     
   substancial dos pobres e vulneráveis.                                        
 • 1.4 até 2030, garantir que todos os homens e mulheres, particularmente os    
   pobres e vulneráveis, tenham direitos iguais aos recursos econômicos, bem    
   como acesso a serviços básicos, propriedade e controle sobre a terra e outras
   formas de propriedade, herança, recursos naturais, novas tecnologias         
   apropriadas e serviços financeiros, incluindo microfinanças.                 
 • 1.5 até 2030, construir a resiliência dos pobres e daqueles em situação de   
   vulnerabilidade, e reduzir a exposição e vulnerabilidade destes a eventos    
   extremos relacionados com o clima e outros choques e desastres econômicos,   
   sociais e ambientais.                                                        
 • 1.a garantir uma mobilização significativa de recursos a partir de uma       
   variedade de fontes, inclusive por meio do reforço da cooperação para o      
   desenvolvimento, de forma a proporcionar meios adequados e previsíveis para  
   que os países em desenvolvimento, em particular os países de menor           
   desenvolvimento relativo, implementem programas e políticas para acabar com a
   pobreza em todas as suas dimensões.                                          
 • 1.b criar marcos políticos sólidos, em níveis nacional, regional e           
   internacional, com base em estratégias de desenvolvimento a favor dos pobres 
   e sensíveis a gênero, para apoiar investimentos acelerados nas ações de      
   erradicação da pobreza. Se quiser, posso fazer o mesmo para os ODS 2, 3, 4...
   em formato resumido ou integral.                                             

Acabar com a fome, alcançar a segurança alimentar e melhoria da nutrição e promover a agricultura sustentável

2.1 Até 2030, acabar com a fome e garantir o acesso de todas as pessoas, em particular os pobres e pessoas em situações vulneráveis, incluindo crianças, a alimentos seguros, nutritivos e suficientes durante todo o ano

2.2 Até 2030, acabar com todas as formas de desnutrição, incluindo atingir, até 2025, as metas acordadas internacionalmente sobre nanismo e caquexia em crianças menores de cinco anos de idade, e atender às necessidades nutricionais dos adolescentes, mulheres grávidas e lactantes e pessoas idosas

2.3 Até 2030, dobrar a produtividade agrícola e a renda dos pequenos produtores de alimentos, particularmente das mulheres, povos indígenas, agricultores familiares, pastores e pescadores, inclusive por meio de acesso seguro e igual à terra, outros recursos produtivos e insumos, conhecimento, serviços financeiros, mercados e oportunidades de agregação de valor e de emprego não agrícola

2.4 Até 2030, garantir sistemas sustentáveis de produção de alimentos e implementar práticas agrícolas resilientes, que aumentem a produtividade e a produção, que ajudem a manter os ecossistemas, que fortaleçam a capacidade de adaptação às mudanças climáticas, às condições meteorológicas extremas, secas, inundações e outros desastres, e que melhorem progressivamente a qualidade da terra e do solo

2.5 Até 2020, manter a diversidade genética de sementes, plantas cultivadas, animais de criação e domesticados e suas respectivas espécies selvagens, inclusive por meio de bancos de sementes e plantas diversificados e bem geridos em nível nacional, regional e internacional, e garantir o acesso e a repartição justa e equitativa dos benefícios decorrentes da utilização dos recursos genéticos e conhecimentos tradicionais associados, como acordado internacionalmente

2.a Aumentar o investimento, inclusive via o reforço da cooperação internacional, em infraestrutura rural, pesquisa e extensão de serviços agrícolas, desenvolvimento de tecnologia, e os bancos de genes de plantas e animais, para aumentar a capacidade de produção agrícola nos países em desenvolvimento, em particular nos países menos desenvolvidos

2.b Corrigir e prevenir as restrições ao comércio e distorções nos mercados agrícolas mundiais, incluindo a eliminação paralela de todas as formas de subsídios à exportação e todas as medidas de exportação com efeito equivalente, de acordo com o mandato da Rodada de Desenvolvimento de Doha

2.c Adotar medidas para garantir o funcionamento adequado dos mercados de commodities de alimentos e seus derivados, e facilitar o acesso oportuno à informação de mercado, inclusive sobre as reservas de alimentos, a fim de ajudar a limitar a volatilidade extrema dos preços dos alimentos

Agenda 2030 no Brasil
Goveno Federal
Existe a Comissão Nacional para os Objetivos de Desenvolvimento Sustentável - CNODS, vinculada à Secretaria-Executiva da Secretaria-Geral da Presidência da República, com a participação de várias instituições como IPEA, IBGE, Fiocruz, Itaipu Binacional, Ministério do Planejamento e Orçamento e Ministério das Relações Exteriores.
O Governo Federal publico em 2024 o Relatório Nacional Voluntário, não publicou em 2025 o relatório de 2026 está em elaboração. O governo Federal se declara comprometido  com a ONU nos seus programas de sustentabilidade.
Segundo o governo federal a lei 14.802/2024 que é o PPA 2024 (Plano Plurianual 2024 - 2027) é um marco na retomada da capacidade de planejamento das açẽs publicas federais, este plano tem a Agenda 2030 como um importante referêncial.
O RNV destaca ainda as ações de planejamento que corroboram com a Agenda 2030: Reformas estruturais e politica de inivação ; o Novo PAC (Programa de Aceleração do Crescimento); Nova Industria Brasil; Politica de Migração e Adaptação Climática; Programa de Aquisição de Alimentos; Ações de Defesa a Democracia; Programas de Incentivo a Cultura;Açoes multidimencionais para a promoção do trabalho decente e dignidade;Programas de insentivo ao esporte e ; diversas comucões com o proposito de incentivar a participação social: Comissão Nacional de População e Desenvolvimento (CNPD), Conselho Nacional dos Direitos da Pessoa Idosa (CNDI), Conselho Nacional de Política Indigenista (CNPI), comitês voltados para os direitos de grupos LGBT-QIA+, crianças e adolescentes, bem como conselhos dedicados à ciência, tecnologia e inovação (CCT), segurança alimentar e nutricional (CONSEA, CAISAN) e meio ambiente (CONAMA, Comitê Interministerial de Mudança do Clima).
Programas
Bolsa Família
Programa nacional de reforma agrária
Programa Federal de Ações Afirmativas

Neste panorama em relação as 169 metas da Agenda a situação se resume em
14 Atingidas
35 Eem evolução
26 Sem Evolução
23 Involução
71 Sem indicadores.
Em relação ao ODS 1 Erradicação da Pobreza o relatório declara:
Houve poucos avanços na redução da pobreza extrema nos últimos dez anos.
A pandemia teve um impacto negativo notável no ODS 1. Com uma breve que-
da da pobreza extrema em 2020 por causa do Auxílio Emergencial, houve
um aumento significativo em 2021, após a interrupção desse auxílio. A re-
cuperação do mercado de trabalho e a retomada das transferências sociais
contribuíram para uma nova queda em 2022. No entanto, para que as metas
do ODS 1 sejam atingidas até 2023, o Brasil ainda precisa reduzir anualmente a taxa
de pobreza extrema em cerca de 0,4 ponto percentual (p.p.). As disparidades na in-
cidência da pobreza persistem entre diferentes regiões e grupos. As taxas mais altas
estão nas regiões norte e nordeste, entre a população preta e parda e nas áreas ru-
rais. A redução da pobreza multidimensional tem sido mais positiva do que a redução
da pobreza monetária, embora haja desafios na obtenção de dados regulares para
monitoramento. Programas de transferência de renda, como o Bolsa Família, desem-
penham um papel crucial na redução da pobreza. A nova fase do Bolsa Família, ini-
ciada em 2023, amplia a cobertura e os valores transferidos, com foco na proteção
à primeira infância, representando uma iniciativa inspiradora para outros países em
desenvolvimento na luta contra a pobreza.

Em relação a Erradicar a Fome
O Brasil enfrenta desafios significativos em relação ao ODS 2. Embora a
fome seja uma questão antiga no país, o conceito de Segurança Alimentar e
Nutricional surgiu nos anos 1980. Desde então, foram desenvolvidas políti-
cas e programas para combater a fome e promover a segurança alimentar ao
mesmo tempo, como o Programa Fome Zero, lançado em 2003. Com isso,
entre 2004 e 2013, houve progresso significativo na redução da fome e o
Brasil, em 2014, saiu do Mapa da Fome da Organização das Nações Unidas (ONU)
pela primeira vez. No entanto, os anos seguintes foram marcados por uma crise
econômica e política, que resultou no aumento da insegurança alimentar, agravado
pela pandemia. Os dados mais recentes mostram que 9,4% dos domicílios do Brasil
estavam em insegurança alimentar moderada ou grave em 2023, indicando uma
redução de 3,3 p.p. em relação a 2017-2018, mas 1,6 p.p. acima de 2013. Portanto,
é necessário ainda um esforço conjunto do governo, da sociedade civil e do setor
privado para promover políticas e programas que garantam o acesso a alimentos
nutritivos e suficientes para toda a população, especialmente a mais vulnerável. Por
essa razão, o PPA 2024-2027 enfatiza o enfrentamento da insegurança alimentar
e da pobreza como prioridades, além do fortalecimento da agricultura familiar e do
agronegócio sustentável.
https://odsbrasil.gov.br/objetivo1/indicador111

Sociedade Civil
Na sociedade civil existe o "Grupo de Trabalho da Sociedade Civil para a Agenda 2030 de Desenvolvimento Sustentável"  formada por 50 entidade [[https://gtagenda2030.org.br/quem-faz/][Lista de entidades do GT Agenda 2030]]. Este grupo publica anualmente um relatório de acompanhamento das
https://gtagenda2030.org.br/quem-faz/e estagnação para a meta 1.b
https://gtagenda2030.org.br/quem-faz/agenda2030.org.br/wp-content/uploads/2025/09/relatorio-luz-2025.pdf

Pela analis do relatório Luz 2025 as Metas 1.1,1.2,1.3 as ações foram consideradas satisfatórias, em relação as metas 1.4,1.5,1.a houve retrocesso, para a meta 1.b o relatório indica estagnação.

Para as metas do ODS 2 o Relatório classifica as ações para as metas 2.1,2.3 como satisfatórias , para as metas 2.2,2.5,2.a insuficientes. Para a meta 2.b as acóes foram consideradas estagnada e as metas 2.4 e 2.c foram consideradas ameaçadas.

PARANA
Estratógia Paraná de Olho nos ODS
https://www.sgdes.pr.gov.br/Pagina/BI-ODS
CURITIBA
Gogar
agenda 2030 no município de Curitiba
Curitiba implementa a Agenda 2030 da ONU por meio da Lei Municipal 15.538/19, que institui os Objetivos de Desenvolvimento Sustentável (ODS) como diretriz permanente de políticas públicas na cidade, independentemente de alternâncias de gestão.Marcos e Governança LocalMarco Legal: Aprovada em 2019, a regulamentação transformou os 17 ODS e as 169 metas da ONU em obrigações de planejamento estratégico municipal.Parcerias Estratégicas: Conta com capacitação de servidores em cooperação com o Centro Internacional de Formação de Autoridades e Líderes (Cifal Curitiba).Transparência e Indicadores: A Prefeitura monitora o avanço das metas locais por meio de um portal de monitoramento gerido pelo Instituto Municipal de Administração Pública ([IMAP](https://www.curitiba.pr.gov.br/noticias/site-revela-progresso-de-curitiba-no- cumprimento-dos-objetivos-do-desenvolvimento-sustentavel/63105)).Principais Eixos de Atuação na CidadeInovação e Cidades Inteligentes: Projetos urbanísticos alinhados ao conceito de sustentabilidade e transição verde, destacados em eventos globais como o Smart City Expo.Desenvolvimento Urbano: Integração de diretrizes de mobilidade e habitação com foco em reduzir desigualdades sociais e otimizar recursos ambientais.Resiliência Climática: Práticas voltadas à eficiência energética, gestão de resíduos e expansão de áreas verdes urbanas.
\end{document}
